\documentclass{article}
\usepackage{amsmath,amssymb,hyperref}
\title{Losing episode identity can square the calibration cost}
\begin{document}
\maketitle

\paragraph{Result.}
For the fixed-cap menu with \(d\) possible bundles, let \(\alpha\)
be the fraction of calibration rows that truly link the necessary
coordinates to one episode. The worst-case number of recorded
rows needed to learn the entire menu is of order
\(d^2\log(d)/\alpha^2\) if those links are hidden, and
\(d^2\log(d)/\alpha\) when they are reliably marked. These rates
are for independent, outcome-independent linking at fixed
familywise error and the fixed exposure-law family below.
They count assembled rows, not the number of original
source episodes consumed or the cost of knowing a link.
Even an honest link flag is insufficient if the linkage process
can select records by the missing joint outcome.

\paragraph{Data and law.}
Use the active-window law from the
\href{../Risk_Bits_Before_The_Query.tex}{risk menu}
with density \(1+\epsilon\sum_j\sigma_j\chi_{S_j}\)
against the uniform law on \(n\) bits, where
\(\epsilon=1/(2d)\), \(|S_j|=m\ge2\), and every candidate
includes the earlier committed variables. In this calibration
study all source episodes are from the active regime.
The sign menu still answers the original joint-cap problem
if the operational window probability is supplied separately;
this does not estimate that probability. The objective is all
signs, not just the mean number of mistakes.

A recorded row first picks a partition \(P\) of the coordinate
labels, independently of outcomes. Each block is filled by
the corresponding coordinates of a distinct independent original
draw from the one fixed law. No original draw is reused across
recorded rows. A tag tells whether a given \(S_j\) lies in one
block; in the hidden case the observer only sees the assembled
bits. Define \(\alpha_j=P(S_j\text{ lies in one block})\) and
assume every \(\alpha_j\ge\alpha>0\). A matcher score, coarsened
timestamp or fraction linked on average is not this assurance.

\paragraph{What the assembled row measures.}
If \(S_j\) lies in one block, its parity expectation is
\(\epsilon\sigma_j\). If it crosses blocks, its parity splits
into the product of characters of nonempty proper subsets
of \(S_j\). Each mean is zero by the original low-order
marginals, and the factors are on independent draws.
Therefore the hidden-row expectation is exactly
\(\alpha_j\epsilon\sigma_j\).
If all sets share a nonempty earlier set, at most one block
can contain any candidate. In that case the \emph{whole}
conditional assembled density is just the constant term plus
the old characters supported in that one block; averaging retains
the same linear density with coefficients
\(\alpha_j\epsilon\sigma_j\). Outside this common-anchor case,
cross-products can appear in the whole density. The individual
parity identity still holds.

\paragraph{Sharp data rates.}
Without a tag, the sign of the sample parity mean uses a
margin of at least \(\alpha\epsilon\). Hoeffding and a union
bound give sufficient size
\(8d^2\alpha^{-2}\log(d/\delta)\).
With a tag, score parity only on rows with the
query correctly linked and set its score to zero otherwise.
Its law has conditional mean \(\epsilon\sigma_j\) on
the informative rows, occurring with probability \(\alpha_j\).
The signed Chernoff bound is
\((1-\alpha_j+\alpha_j\sqrt{1-\epsilon^2})^N\).
Thus \(8d^2\alpha^{-1}\log(d/\delta)\) suffices.
Neither sign rule needs the exact linkage rate; fixed
budgeting uses the stated lower bound.

For the lower bounds choose the partition to keep every
coordinate together with probability \(\alpha\) and otherwise
separate all coordinates. With singleton blocks the row is
uniform under every menu. Under a hidden partition, the
observed density is \(1+\alpha\epsilon\sum_j\sigma_j\chi_{S_j}\).
An alternative that flips one sign differs from the all-positive
reference by \(2\alpha\epsilon\) in each relative density cell,
with denominator at least \(1/2\). Its KL is at most its
chi-square divergence, hence at most \(2\alpha^2/d^2\).
With a supplied tag the conditional unlinked component has
no divergence and the linked component's KL is at most
\(2/d^2\), so the joint KL is at most \(2\alpha/d^2\).
Use the uniform prior on the reference and the \(d\)
single-flip alternatives. Average KL to the reference
bounds mutual information; independence makes it scale
with the number of recorded rows. Fano then requires
\[
 \frac{d^2}{2\alpha^b}
 \big[(1-\delta)\log(d+1)-\log2\big]
\]
rows, with \(b=2\) for hidden and \(b=1\) for tagged linkage
when the bracket is positive. This matches the stated rates at
fixed error up to constants. These lower examples use the
same partitions and candidate menu in the two information
interfaces.

\paragraph{When tags do not help.}
Take two menu laws with opposite exchanged signs at \(S_1,S_2\)
and all other signs the same. Apply the \emph{same} outcome-dependent
retention rule in both: keep the whole source episode exactly
when \(\chi_{S_1}=\chi_{S_2}\). On that event the two original
densities coincide. The event has probability \(1/2\);
the symmetric difference \(S_1\mathbin{\triangle}S_2\) excludes
the common earlier set and is not a menu character, so its
mean is zero. On rejected rows, form singleton blocks from
independent rejected source trials. The two laws had identical
singleton marginals and identical retained subdistributions,
so their rejected singleton marginals also match.
The entire tagged data law is now the same in both worlds,
while two decisions at the original cap differ. No amount of
this data identifies them. A correct retention indicator
does not say why those rows were retained.

\paragraph{Checks and relation to earlier work.}
Null and positive cases check 168,896 exact assembled-law
cells across all 203 six-coordinate partitions, ten one-flip
chi-square controls, a disjoint-block case with a new
cross-character, and the single shared censor rule above.
The classical statistical concern that linkage errors can
distort downstream inference appears in
\href{https://www150.statcan.gc.ca/n1/en/catalogue/12-001-X199300114476}{Scheuren
and Winkler}. Here the distinction on the same fixed menu is
observable linkage versus latent linkage versus selection on
unseen joint outcomes. The experiment does not validate any
actual link source, independence premise, or R4 acceptance.
\end{document}
